\clearpage
\section{Statistical methods and interpretation}\label{app:statistics}

\subsection{Bootstrapping}\label{app:bootstrap}
To estimate how repeat measurements affect a fitted result, we use bootstrapping. We repeatedly draw measurements from the available dataset \emph{with replacement}. A measurement can therefore be selected more than once, while another may be left out. Each resample has the same size as the original dataset. Repeating the analysis for these resamples gives a distribution of fitted values, whose spread describes the repeat variability.

We resample complete retained scans. If a file contains $n$ scans $S_1(t),\ldots,S_n(t)$ at the same delays, we draw $n$ indices independently and uniformly from $\{1,\ldots,n\}$ for each resample $r$. Its mean is
\[
J_1^{(r)},\ldots,J_n^{(r)},\qquad
\overline S^{*(r)}(t)=\frac1n\sum_{j=1}^{n}S_{J_j^{(r)}}(t).
\]
We fit this mean using the same function, interval, and procedure as for the original data. Repeating the calculation $R=300$ times gives lifetime estimates $\tau^{*(1)},\ldots,\tau^{*(R)}$. The value reported in the results is still the fit to the original mean; the resampled fits are used to estimate its uncertainty.

For quantities involving several files, scans are resampled separately within each file. Each replicate then repeats all relevant fits and derived calculations. In particular, slow lifetimes are converted to rates and regressed against the unchanged $q^2$ coordinates, treated as exact to obtain a bootstrap diffusion estimate. The paired fast/slow diagnostic uses the same resampled scans for both components, preserving their correlation.

Resampling whole scans preserves relationships between delay points within a scan. Treating all individual delay points as independent would discard that structure. The method nevertheless assumes that the retained scans adequately represent repeat measurements under comparable conditions; it cannot recover sources of variability absent from those scans.

\subsection{Repeat-scan bootstrap intervals}\label{app:bootstrap-intervals}
A \emph{repeat-scan bootstrap interval} summarizes the spread of the estimates obtained by the procedure above. The conditional 95\% percentile interval for a parameter $\theta$ is
\[
\mathrm{CI}_{95\%}(\theta)=
\left[Q_{0.025}(\theta^*),\;Q_{0.975}(\theta^*)\right],
\]
where $Q_p$ denotes the $p$ quantile of the bootstrap estimates. Its endpoints are the 2.5th and 97.5th percentiles, enclosing the middle 95\% of those estimates. The interval need not be symmetric about the original estimate.

The quantiles are calculated by sorting the $N$ bootstrap estimates as $x_0\leq x_1\leq\cdots\leq x_{N-1}$ and applying NumPy's default linear interpolation. With indices starting at zero and $0<p<1$, define
\[
h=(N-1)p,\qquad j=\lfloor h\rfloor,\qquad a=h-j.
\]
The quantile is the interpolated value
\[
Q_p=(1-a)x_j+a x_{j+1}.
\]
An integer $h$ gives $a=0$ and selects a sorted estimate directly; otherwise the quantile lies between its two neighboring estimates. The implementation uses \texttt{np.percentile} with the requested percentages 2.5 and 97.5.

For the $N=300$ replicates used here, the lower bound is
\[
h=299(0.025)=7.475,
\qquad Q_{0.025}=0.525x_7+0.475x_8,
\]
between the eighth and ninth smallest estimates. The upper bound is
\[
h=299(0.975)=291.525,
\qquad Q_{0.975}=0.475x_{291}+0.525x_{292},
\]
between the 292nd and 293rd smallest estimates. These bounds are taken from the ordered bootstrap distribution, rather than calculated as the fitted value plus or minus a multiple of its standard deviation.

These intervals estimate uncertainty arising from variation among the retained repeated scans, conditional on the chosen model, fit window, and fixed scan-selection masks. The supplied PP/TG delays and grating spacings are assumed exact. The intervals do not include uncertainty in the excluded-scan selection, unrepresented acquisition drift, or an unsuitable physical model. A narrow interval can therefore coexist with systematic residuals. The interval is also different from a conditional covariance standard error written as $\pm u$: the two calculations describe different sources of statistical variability. A 95\% confidence interval is not a statement that 95\% of individual measurements lie inside its bounds, or a guaranteed probability statement about this particular realized interval.

Derived quantities are recalculated for each resample. For example, $k^{*(r)}=1/\tau^{*(r)}$ supplies a bootstrap distribution of rates. Since the reciprocal is decreasing for positive lifetimes, lifetime bounds $[L,U]$ correspond to rate bounds $[1/U,1/L]$. The plotted rate intervals are calculated directly from the resampled rates. Neither symmetric lifetime errors nor their reciprocals should be treated as automatically symmetric rate errors.

\subsection{Coefficient of determination \texorpdfstring{$R^2$}{R-squared}}\label{app:r-squared}
The coefficient of determination $R^2$ describes how much of the variation in the fitted decay rates is accounted for by a straight line. For rates $k_i=1/\tau_{s,i}$, their mean $\overline{k}$, and the line predictions $\widehat{k}_i=b+D_aq_i^2$, it is defined as
\[
R^2=1-\frac{\sum_i(k_i-\widehat{k}_i)^2}
{\sum_i(k_i-\overline{k})^2}.
\]
The numerator measures the squared deviations from the fitted line; the denominator measures the squared deviations from the mean rate. For the unweighted least-squares regressions with a freely fitted intercept used here, $R^2=1$ means that all points lie on the line, whereas $R^2=0$ means that the line predicts no better than the constant mean. This interpretation requires nonzero variation in the rates.

The TG-only value $R^2=0.934$ corresponds to 93.4\% of the rate variation accounted for by the line; including PP gives $R^2=0.161$, or 16.1\%. These are descriptive statistics with unestimated uncertainties, calculated for four and five points, respectively. Their comparison is not a formal statistical test. A high $R^2$ alone does not validate the physical model, quantify parameter uncertainty, or account for model mismatch and systematic residuals. It also does not show that the fitted slow components represent the same microscopic process in every dataset.

\subsection{Reduced chi-squared diagnostics}\label{app:chi-squared}
If a fitted curve $f(x_i;\widehat{\boldsymbol\theta})$ is compared with $N$ observations $y_i$ having standard errors $s_i$, a diagonal reduced chi-squared statistic is
\[
\chi^2_\nu=\frac{1}{\nu}\sum_{i=1}^{N}
\frac{[y_i-f(x_i;\widehat{\boldsymbol\theta})]^2}{s_i^2},
\qquad \nu=N-p,
\]
where $p$ is the number of independently fitted parameters. The nominal denominators use $p=5$ for a biexponential plus constant, $p=7$ for three exponentials plus constant, and $p=2$ for a straight line. The numerical values are descriptive diagnostics; uncertainties on these statistics have not been estimated.

For PP and TG mean traces, $s_i$ is the sample standard deviation across retained scans at delay $t_i$, divided by the square root of the number of retained scans. For rate regressions, $s_i$ is the standard deviation of that rate's 300 whole-scan bootstrap estimates. It is not the half-width of its 95\% percentile interval. The regression retaining borderline scan 4 uses a separate resampling calculation for that scan selection. The supplied delay and grating coordinates remain exact inputs.

The fits minimize unweighted squared residuals, while this statistic divides the residuals by estimated errors afterwards. In addition, delay points within a scan may be correlated, and the errors are estimated from a finite number of repeats. The diagonal sum neglects these correlations. We therefore use $N-p$ as a nominal number of degrees of freedom and do not assign a formal goodness-of-fit probability or rejection threshold. A large value indicates deviations that are large compared with the estimated repeat variability. It must be interpreted together with the residual plots; a value near unity alone does not validate the model.

The single spectrum and IAC records have no supplied pointwise measurement errors or independent noise model, so their reduced chi-squared is marked unavailable. This also applies to their fit-window variants. Setting a common noise variance equal to $\widehat s^2=\sum_i[y_i-f(x_i)]^2/(N-p)$ would force $\chi^2_\nu=1$ by construction and provide no independent check. The fringe-frequency curves are derivatives of the IAC fits, not separate fits to the blue points; no additional reduced chi-squared is assigned to them.

The main-fit diagnostics accompany the figures. The following values cover the additional fit-window and all-scan alternatives quoted in the text. Scan-selection masks remain fixed when varying the fit window. Each rate-line variant uses rate standard deviations from a separate 300-replicate whole-scan bootstrap; the supplied delay and grating coordinates remain fixed. The same diagonal and unweighted-fit limitations apply, and uncertainties on these diagnostic values have not been estimated.
% Generated by analysis/figure_fit_details.py.
\begin{center}\small\begin{tabular}{lrr}\toprule
Additional fit & Diagonal $\chi^2_\nu$ & $\nu$\\\midrule
PP, 0.5--33 ps & 0.141 & 19\\
PP, 0.5--57 ps & 1.44 & 20\\
PP, 0.5--100 ps & 1.89 & 21\\
PP, 0.5--2000 ps & 120 & 59\\
PP, 1--2000 ps & 129 & 54\\
PP, 2--2000 ps & 11.1 & 52\\
TG slow-rate line, start 0.2 ps & 1.32 & 2\\
TG slow-rate line, start 1 ps & 2.46 & 2\\
TG slow-rate line, start 2 ps & 0.168 & 2\\
TG slow-rate line, all scans retained & 0.591 & 2\\
\bottomrule\end{tabular}\end{center}

